\documentclass[11pt]{article}
\usepackage{mathblatt}
\begin{document}
% Kopfzeile Seite 1 ohne Kennung; die Rückseite (Lösungen, für den Lehrer) bekommt sie nach \begleitteil
\blattkopf{Mathematik}{Basisaufgaben}
\einheitenkopf[][Zettel 18]{Basisaufgaben · Zettel 18}
\zweigzeile{je Aufgabe 1 Punkt · Lösungen auf der Rückseite}
\par\vspace{2pt}\noindent Name \underline{\hspace{7cm}}\hfill Datum \underline{\hspace{3cm}}\hfill Punkte \underline{\hspace{1.2cm}} / 18\par\vspace{6pt}
% \small: volle Seite, kurze Aufgaben zweispaltig (v0.6)
\small

\noindent\begin{minipage}[t]{0.485\linewidth}\raggedright
% 1: Bruchteil einer Größe berechnen – brueche-dezimalzahlen-basis-k2-v7 (Original 2018-OS-B1a, ausgedacht: keine Marke)
\begin{aufgabe}{$\frac{2}{3}$ von $45$ min – wie viele Minuten? \leerfeld[min]}
\end{aufgabe}
\end{minipage}\hfill\begin{minipage}[t]{0.485\linewidth}\raggedright
% 2: Mitte zweier Zahlen bestimmen – brueche-dezimalzahlen-basis-k3-v3 (Original 2020-OS-B1f, ausgedacht: keine Marke)
\begin{aufgabe}{Welche Zahl liegt genau in der Mitte zwischen $0{,}7$ und $0{,}8$? \leerfeld}
\end{aufgabe}
\end{minipage}\par

\noindent\begin{minipage}[t]{0.485\linewidth}\raggedright
% 3: Zahl zu Bedingung angeben – rationale-zahlen-basis-k3-v5 (Original 2015-OS-B1b, ausgedacht: keine Marke)
\begin{aufgabe}{Gib eine Zahl an, die zwischen $0{,}6$ und $0{,}7$ liegt. \leerfeld}
\end{aufgabe}
\end{minipage}\hfill\begin{minipage}[t]{0.485\linewidth}\raggedright
% 4: Zehnerpotenzschreibweise umwandeln – potenzen-wurzeln-basis-k3-v10 (Original 2025-OS-B1d, ausgedacht: keine Marke)
\begin{aufgabe}{$0{,}7 = 7 \cdot 10^{\square}$ – welche Hochzahl? \leerfeld}
\end{aufgabe}
\end{minipage}\par

\noindent\begin{minipage}[t]{0.485\linewidth}\raggedright
% 5: Prozentsatz berechnen – prozentrechnung-basis-k3-v3 (Original 2014-OS-B1e, ausgedacht: keine Marke)
\begin{aufgabe}{In einer Lostrommel sind $80$ Lose, $12$ davon gewinnen. Wie viel Prozent sind das? \leerfeld[\%]}
\end{aufgabe}
\end{minipage}\hfill\begin{minipage}[t]{0.485\linewidth}\raggedright
% 6: Größen vergleichen – einheiten-basis-k2-v5 (Original 2026-FOR-B1d, ausgedacht: keine Marke)
\begin{aufgabe}{Vergleiche $2{,}4$ t und $240$ kg. Setze das richtige Zeichen ein: $<$, $=$ oder $>$. $2{,}4$ t \leerfeld $240$ kg}
\end{aufgabe}
\end{minipage}\par

\noindent\begin{minipage}[t]{0.485\linewidth}\raggedright
% 7: Proportionale Zuordnung Dreisatz – zuordnungen-basis-k2-v5 (Original 2023-OS-B1a, ausgedacht: keine Marke)
\begin{aufgabe}{Für $10$ Kekse braucht man $250$ g Mehl. Wie viel Mehl braucht man für $24$ Kekse? \leerfeld[g]}
\end{aufgabe}
\end{minipage}\hfill\begin{minipage}[t]{0.485\linewidth}\raggedright
% 8: Lösung durch Einsetzen prüfen – quadratische-gleichungen-basis-k1-v9 (Original 2025-OS-B1h, ausgedacht: keine Marke)
\begin{aufgabe}{Welcher Wert erfüllt $x \cdot (x + 10) = -21$? \\ \kreuz{$x = 3$} \kreuz{$x = 7$} \kreuz{$x = -7$} \kreuz{$x = -21$}}
\end{aufgabe}
\end{minipage}\par

% 9: Rechteckseite aus Fläche berechnen – flaechen-basis-k3-v2 (Original 2023-OS-B1c, ausgedacht: keine Marke)
\begin{aufgabe}{Ein rechteckiges Spielfeld hat den Flächeninhalt $1800$ m² und ist $30$ m breit. Wie lang ist es? \leerfeld[m]}
\end{aufgabe}

% 10: Portionen aus Gesamtmenge berechnen – einheiten-basis-k3-v2 (Original 2022-OS-B1h, ausgedacht: keine Marke)
\begin{aufgabe}{Ein Bäcker hat $3$ kg Mehl. Für einen Kuchen braucht er $500$ g. Für wie viele Kuchen reicht das Mehl? \leerfeld Kuchen}
\end{aufgabe}

% 11: Punktprobe durchführen – lineare-funktionen-basis-k3-v3 (Original 2023-OS-B1i, ausgedacht: keine Marke)
\begin{aufgabe}{Gegeben ist $f(x) = 4x - 3$. Welcher Punkt liegt nicht auf der Geraden? Kreuze an. \\ \kreuz{$(-1|-7)$} \kreuz{$(2|8)$} \kreuz{$(3|9)$}}
\end{aufgabe}

% 12: Lage eines Punktes zu den Achsen erkennen – symmetrie-abbildungen-basis-k2-v3 (Original 2020-OS-B1b, ausgedacht: keine Marke)
\begin{aufgabe}{Genau einer der vier Punkte liegt auf der x-Achse. Welcher? \\ \kreuz{$K(-8|0)$} \kreuz{$L(0|-8)$} \kreuz{$M(-8|-8)$} \kreuz{$N(8|8)$}}
\end{aufgabe}

% 13: Eigenschaft einer Figur zuordnen – winkel-dreiecke-basis-k5-v5 (Original 2026-FOR-B1c, ausgedacht: keine Marke)
\begin{aufgabe}{Kreuze die richtige Ergänzung an: „In jedem gleichseitigen Dreieck …“\\ \kreuz{… gibt es einen rechten Winkel.}\\ \kreuz{… ist jeder Winkel $60^\circ$ groß.}\\ \kreuz{… ist ein Winkel stumpf.}}
\end{aufgabe}

% 14: Kantenzahl eines Körpers angeben – koerper-basis-k2-v3 (Original 2019-OS-B1f, ausgedacht: keine Marke)
\begin{aufgabe}{Ein Turmdach hat die Form einer Pyramide mit fünfeckiger Grundfläche. Gib an, wie viele Kanten diese Pyramide hat. \leerfeld Kanten}
\end{aufgabe}

% 15: Arithmetisches Mittel berechnen – daten-basis-k1-v5 (Original 2021-OS-B1f, ausgedacht: keine Marke)
\begin{aufgabe}{An der Bushaltestelle wartet Ali $5$; $9$; $7$; $11$ min. Berechne das arithmetische Mittel (Wartezeit). \leerfeld[min]}
\end{aufgabe}

% 16: Zufallsgerät zu Wahrscheinlichkeit entwerfen – wahrscheinlichkeit-basis-k2-v4 (Original 2019-OS-B1g, ausgedacht: keine Marke)
\begin{aufgabe}{In einer Tüte sind $5$ grüne Bonbons und noch andere. Die Wahrscheinlichkeit für Grün soll $\frac{1}{3}$ sein. Wie viele andere Bonbons müssen in der Tüte sein? \leerfeld Bonbons}
\end{aufgabe}

% 17: Bruchteil einer Fläche bestimmen – brueche-dezimalzahlen-basis-k1-v13 (Original 2026-FOR-B1b, ausgedacht: keine Marke)
\begin{aufgabe}{\begin{minipage}[t]{0.6\linewidth}Ein Beet ist in gleich große Streifen geteilt; auf den grauen wachsen Kräuter. Wie viel Prozent der Fläche sind grau? \leerfeld[\%]\end{minipage}\hfill\begin{minipage}[t]{0.37\linewidth}\vspace{0pt}
\bruchrechteck[4]{4}{5}
\end{minipage}}
\end{aufgabe}

% 18: Winkel an geschnittenen Parallelen bestimmen – winkel-dreiecke-basis-k3-v3 (Original 2020-OS-B1g, ausgedacht: keine Marke)
\begin{aufgabe}{\begin{minipage}[t]{0.6\linewidth}Die beiden waagerechten Geraden sind parallel. Wie groß ist $\alpha$? $\alpha =$ \leerfeld °\end{minipage}\hfill\begin{minipage}[t]{0.37\linewidth}\vspace{0pt}
\parallelenpaar{70}{70^\circ}{}{\alpha}{}
\end{minipage}}
\end{aufgabe}

\begleitteil
\blattkopf{Basisaufgaben · Zettel 18}{BAS-Z18 · Lösungen}
\erg{1}{$45 : 3 \cdot 2 = 30$ min \hfill {\footnotesize\textit{Bruchteil einer Größe berechnen · 2018}}}
\erg{2}{$(0{,}7 + 0{,}8) : 2 = 0{,}75$ \hfill {\footnotesize\textit{Mitte zweier Zahlen bestimmen · 2020}}}
\erg{3}{z. B. $0{,}65$; richtig ist jede Zahl größer als $0{,}6$ und kleiner als $0{,}7$ \hfill {\footnotesize\textit{Zahl zu Bedingung angeben · 2015}}}
\erg{4}{$-1$ – das Komma wandert eine Stelle nach rechts \hfill {\footnotesize\textit{Zehnerpotenzschreibweise umwandeln · 2025}}}
\erg{5}{$12 : 80 = 0{,}15 = 15\,\%$ \hfill {\footnotesize\textit{Prozentsatz berechnen · 2014}}}
\erg{6}{$2{,}4$ t $= 2400$ kg, also $2{,}4$ t $>$ $240$ kg \hfill {\footnotesize\textit{Größen vergleichen · 2026}}}
\erg{7}{Für $1$: $250 : 10 = 25$; $25 \cdot 24 = 600$ g \hfill {\footnotesize\textit{Proportionale Zuordnung Dreisatz · 2023}}}
\erg{8}{$x = -7$, denn $(-7) \cdot (-7 + 10) = (-7) \cdot 3 = -21$ \hfill {\footnotesize\textit{Lösung durch Einsetzen prüfen · 2025}}}
\erg{9}{$1800 : 30 = 60$ m \hfill {\footnotesize\textit{Rechteckseite aus Fläche berechnen · 2023}}}
\erg{10}{$3$ kg $= 3000$ g; $3000 : 500 = 6$ \hfill {\footnotesize\textit{Portionen aus Gesamtmenge berechnen · 2022}}}
\erg{11}{$(2|8)$, denn $f(2) = 5$ \hfill {\footnotesize\textit{Punktprobe durchführen · 2023}}}
\erg{12}{$K(-8|0)$ \hfill {\footnotesize\textit{Lage eines Punktes zu den Achsen erkennen · 2020}}}
\erg{13}{… ist jeder Winkel $60^\circ$ groß. \hfill {\footnotesize\textit{Eigenschaft einer Figur zuordnen · 2026}}}
\erg{14}{$5 + 5 = 10$ Kanten ($5$ Grundkanten, $5$ Seitenkanten) \hfill {\footnotesize\textit{Kantenzahl eines Körpers angeben · 2019}}}
\erg{15}{$32 : 4 = 8$ min \hfill {\footnotesize\textit{Arithmetisches Mittel berechnen · 2021}}}
\erg{16}{$5$ grüne sind $\frac{1}{3}$ von $15$, also $15 - 5 = 10$ andere \hfill {\footnotesize\textit{Zufallsgerät zu Wahrscheinlichkeit entwerfen · 2019}}}
\erg{17}{$4$ von $5$ gleichen Teilen: $\frac{4}{5} = 80\,\%$ \hfill {\footnotesize\textit{Bruchteil einer Fläche bestimmen · 2026}}}
\erg{18}{$\alpha = 70^\circ$ (Stufen- oder Wechselwinkel) \hfill {\footnotesize\textit{Winkel an geschnittenen Parallelen bestimmen · 2020}}}
\end{document}
